\section{Обзор существующих решений}
\label{sec:Section2} \index{Section2}

Системы анализа правил фильтрации развивались вместе с сетевыми устройствами. На смену устройствам с простыми пакетными фильтрами приходили устройства, отслеживающие состояние соединений (stateful), а затем появились NGFW с инспекцией приложений и другими расширенными функциями. Менялись и сами правила фильтрации: они уже не ограничивались стандартной пятёркой полей (IP-адрес источника, IP-адрес получателя, порт источника, порт получателя, протокол), а получили более широкий набор квалификаторов; появилась семантика цепочек, изменился подход к организации списка правил. Более ранние работы, ориентированные на простые пакетные фильтры, оказались неприменимы к конфигурациям современных устройств. Последующие подходы перешли от попарного анализа правил к символьным методам и от пятёрки полей к расширенному набору квалификаторов; тем не менее каждый такой переход рассчитывался под конкретную платформу и определённый формат правил и не пытался учесть всё многообразие возможных конфигураций.

Для построения собственной системы в настоящем разделе рассматриваются модели правил фильтрации, уже предложенные в литературе, обсуждается сложность реальных конфигураций, систематизируются типы ошибок и аномалий, выделяемых разными авторами, и анализируются предложенные методы анализа правил --- от попарных и геометрических до символьных, логических, формально верифицированных и распределённых. Отдельно рассматриваются методы оптимизации политик: переупорядочивание правил и проектирование политики с нуля.

\subsection{Правила фильтрации}
\label{sec:Section2:rules}

В наиболее общем виде правило фильтрации представляет собой пару $\langle P, a\rangle$, где $P$ --- предикат над набором атрибутов сетевого пакета (адрес отправителя, адрес получателя, протокол, порты и иные поля) и контекста установленной сессии (состояние соединения, идентификаторы приложений, URL-категории и т.п.), а $a$ --- действие, применяемое к пакету/сессии в случае истинности предиката.

Правила фильтрации организованы в упорядоченные наборы (списки), которые применяются к каждому проходящему пакету последовательно: первое сработавшее правило определяет результат обработки пакета. Из-за простоты такого определения на правило фильтрации не накладываются практически никакие ограничения, и производители сетевых устройств реализуют данный функционал по-своему: различается состав квалификаторов и способы их задания, множество возможных действий и принципы организации списков правил.

\subsubsection*{Примеры правил фильтрации}

Для примера рассмотрим типовую политику фильтрации трафика на разных сетевых устройствах (на разных типах систем фильтрации трафика). Возьмём следующую задачу: \emph{разрешить HTTPS-трафик из внутренней сети 192.168.1.0/24 к веб-серверу 10.0.0.10 и заблокировать всё остальное}.

В качестве классических правил фильтрации рассмотрим списки контроля доступа (ACL) на маршрутизаторах и коммутаторах \textbf{CISCO IOS}. Данная политика записывается в виде двух строк access-list'а, привязываемого к интерфейсу:

\begin{verbatim}
access-list 101 permit tcp 192.168.1.0 0.0.0.255 host 10.0.0.10 eq 443
access-list 101 deny   ip any any
interface GigabitEthernet0/1
 ip access-group 101 in
\end{verbatim}

Первая команда разрешает TCP-трафик из подсети 192.168.1.0/24 в подсеть 10.0.0.10/32 на порт 443. Вторая команда запрещает всё остальное. При этом данный набор правил действует только на входящий трафик интерфейса GigabitEthernet0/1.

На устройствах c Linux та же политика записывается с помощью команд \texttt{iptables}:

\begin{verbatim}
iptables -A FORWARD -s 192.168.1.0/24 -d 10.0.0.10 \
         -p tcp --dport 443 -j ACCEPT
iptables -A FORWARD -j DROP
\end{verbatim}

Каждая команда добавляет (\texttt{-A}, \emph{append}) одно правило в конец цепочки \texttt{FORWARD} --- предопределённой цепочки таблицы \texttt{filter}, через которую проходит транзитный трафик. Первое правило задаёт условие сопоставления: адрес отправителя из подсети \texttt{192.168.1.0/24} (\texttt{-s}), адрес получателя \texttt{10.0.0.10} (\texttt{-d}), протокол TCP (\texttt{-p tcp}) и порт назначения \texttt{443} (\texttt{--dport}). При выполнении всех условий применяется действие \texttt{-j ACCEPT}, разрешающее пакет. Второе правило не содержит условий (а значит, сопоставляется любому пакету) и применяет действие \texttt{DROP} --- блокировка без уведомления отправителя.

В отличие от Cisco IOS в iptables действие правила задаётся не отдельным ключевым словом (\texttt{permit}/\texttt{deny}), а так называемой \emph{целью} (target), указываемой после флага \texttt{-j}. Целью может быть встроенное действие --- \texttt{ACCEPT} (разрешить), \texttt{DROP} (блокировка без уведомления отправителя), \texttt{REJECT} (блокировка с ICMP-уведомлением), \texttt{LOG} (записать пакет в журнал), \texttt{RETURN} (возврат к предыдущей цепочке) --- либо имя пользовательской цепочки правил, тогда происходит переход на указанную цепочку, а при достижении её конца возвращается обратно. Помимо предопределённых цепочек \texttt{INPUT}, \texttt{OUTPUT} и \texttt{FORWARD}, через которые автоматически проходит соответствующий трафик, администратор может создавать собственные цепочки и использовать их как переиспользуемые блоки правил, что делает iptables существенно более выразительным средством, чем плоский ACL.

В \textbf{NGFW} правило стараются задавать над именованными объектами, зонами и идентификаторами приложений уровня L7, а не над портами. В качестве примера приведём фрагмент конфигурации FortiGate, который реализует политику, описанную выше:

\begin{verbatim}
config firewall address
    edit "Internal-Net"
        set subnet 192.168.1.0 255.255.255.0
    next
    edit "Web-Server"
        set subnet 10.0.0.10 255.255.255.255
    next
end

config application list
    edit "Allow-HTTPS-only"
        config entries
            edit 1
                set application 40568
            next
        end
        set other-application-action block
    next
end

config firewall policy
    edit 1
        set name "Allow-HTTPS"
        set srcintf "internal"
        set dstintf "wan1"
        set srcaddr "Internal-Net"
        set dstaddr "Web-Server"
        set service "ALL"
        set action accept
        set application-list "Allow-HTTPS-only"
        set ssl-ssh-profile "deep-inspection"
        set utm-status enable
    next
    edit 2
        set name "Deny-Default"
        set srcintf "any"
        set dstintf "any"
        set srcaddr "all"
        set dstaddr "all"
        set service "ALL"
        set action deny
    next
end
\end{verbatim}

Конфигурация FortiGate представляет собой иерархию секций \texttt{config ... end}, внутри которых отдельные объекты задаются блоками \texttt{edit \emph{name} ... next}. В секции \texttt{firewall address} создаются именованные сетевые объекты. В секции \texttt{application list} описывается профиль контроля приложений \texttt{Allow-HTTPS-only}, разрешающий лишь приложение с числовым идентификатором сигнатуры HTTPS (40568) и блокирующий любые другие приложения (\texttt{set other-application-action block}). В секции \texttt{firewall policy} задаются два правила с приоритетами 1 и 2. Каждое правило связывает пару интерфейсов (\texttt{srcintf}/\texttt{dstintf}) и пару адресных объектов (\texttt{srcaddr}/\texttt{dstaddr}). В отличие от Cisco IOS правило не привязывается к одному интерфейсу с направлением, а явно указывает входящий и исходящий интерфейсы. У других производителей встречается привязка правил к зонам безопасности (zone-based). Под зоной безопасности понимается группа интерфейсов, которые объединены в единую логическую сущность.

Принципиальное отличие от классических ACL в том, что сопоставление трафика как HTTPS выполняется не по номеру TCP-порта, а методами глубокой инспекции пакетов (DPI). Поле \texttt{service} в правиле отвечает только за фильтрацию по L4 модели OSI (протокол и порт). Идентификация приложения задаётся с помощью профиля контроля приложений (\texttt{application-list}), в котором FortiGate определяет HTTPS даже при использовании нестандартного порта. Поскольку HTTPS-трафик представляет собой связку SSL и HTTP, то для определения HTTP внутри SSL-сессии необходимо использовать профиль SSL/TLS-инспекции (\texttt{ssl-ssh-profile "deep-inspection"}). Опция \texttt{utm-status enable} включает применение UTM-профилей (контроль приложений, IPS, антивирус и пр.) к данному правилу.

Уже на этих простейших примерах видны основные различия и эволюция в подходах к описанию правил фильтрации сетевого трафика. Во-первых, меняется способ привязки правила внутри политик: в Cisco IOS список правил жёстко связан с конкретным интерфейсом и направлением трафика, в iptables --- с цепочкой обработки в подсистеме ядра, а в FortiGate правило формулируется уже над парой интерфейсов или зон безопасности и тем самым отделяется от физической топологии. Во-вторых, меняется форма задания сетевых квалифакторов: первые два примера опираются на примитивы --- IP-адреса, маски и номера портов, --- тогда как современные решения работают с именованными объектами и их группами (\texttt{Internal-Net}, \texttt{Web-Server}, \texttt{Allow-HTTPS-only}), что позволяет описывать политику в терминах организационных сущностей, а не конкретных значений.

Однако отличия могут наблюдаться и, казалось бы, на близких платформах с одинаковым видом правил: имена объектов чувствительны к регистру, использование разного вида объектов, различные нотации для масок и диапазонов, синонимы протоколов, а также различные способы задания действий и т.д. Подобные отличия должны устраняться на этапе \emph{нормализации} правил --- приведении исходной конфигурации к единой внутренней модели правил и объектов, с которой работает система анализа правил. При нахождении новых расхождений может потребоваться не только адаптировать парсер для новой платформы, но и расширить внутреннюю модель правил, что приводит к значительным изменениям в самой системе анализа правил.

\subsubsection*{Состав квалификаторов}

В простых пакетных фильтрах квалификаторы образуют пятёрку полей сетевого и транспортного уровней модели OSI: адрес и порт источника, адрес и порт получателя, протокол. Эта же пятёрка лежит в основе большинства формальных моделей анализа правил, рассматриваемых далее в настоящем разделе. Однако современные межсетевые экраны давно вышли за её пределы. Помимо классической пятёрки в правилах могут участвовать:

\begin{itemize}
    \item название сетевых интерфейсов с учётом направления и зон безопасности;
    \item состояния соединения;
    \item идентификаторы приложений;
    \item URL-категории и доменные имена;
    \item идентификаторы пользователей;
    \item временные интервалы (расписания);
    \item и др.
\end{itemize}

При этом значения квалификаторов могут задаваться как примитивами (например, одиночный IP-адрес, диапазон портов), так и ссылками на именованные объекты и группы. На современных устройствах настройка политики с помощью именованных объектов является хорошей практикой, позволяющей управлять политикой с привязкой к абстрактным организационным сущностям, а не конкретным IP-адресам и портам. Система анализа правил должна уметь работать как с примитивами, так и с именованными объектами и их группами.

\subsection{Действия в правилах}

Самая простая система действий, которая используется в пакетных фильтрах --- бинарная: разрешить и запретить (\emph{allow} / \emph{deny}) прохождение трафика; именно она лежит в основе большинства моделей анализа правил~\cite{al2003firewall,yuan2006fireman,liu2008firewall}. Однако реально производители сетевых устройств предлагают существенно более широкий набор действий:

\begin{itemize}
    \item Различные виды блокировки: \emph{drop} --- блокировка без уведомления отправителя; \emph{reject} --- блокировка с отправкой ICMP-уведомления или TCP RST;
    \item Действия с побочным эффектом, не определяющие судьбу пакета: \emph{log} (запись пакета в журнал), \emph{count} (увеличение счётчика, привязанного к правилу), \emph{mark} (установка метки, читаемой последующими правилами и подсистемой маршрутизации), \emph{queue} (передача пакета пользовательскому процессу для внешнего разбора);
    \item Управление механизмом цепочек в iptables/nftables: \emph{jump} --- переход в пользовательскую цепочку правил и \emph{return} --- возврат из неё;
    \item Специфичные действия для NGFW: \emph{enforce} (требование пройти дополнительные проверки --- IPS, антивирус, DLP), \emph{tunnel} (перенаправление в шифрованный канал) и другие;
    \item Комбинированные действия: \emph{accept-and-log} / \emph{deny-and-log}, упомянутые в работе~\cite{gouda2004firewall}.
\end{itemize}

Так, для пакета, попавшего под правило с действием \emph{enforce}, последует проверка движками безопасности, окончательное решение (\emph{allow} или \emph{deny}) принимается уже не списком правил, а внешним движком безопасности; поэтому с точки зрения чисто фильтрационной модели такое правило разрешает трафик.

Для анализа конфликтов и поиска аномалий нет особого смысла различать разные действия, приводящие к блокировке, так как они одинаково конфликтуют с действиями, разрешающими трафик, и наоборот. Поэтому достаточно различать только два класса --- \emph{allow} и \emph{deny}, а другие действия можно сводить к одному из этих двух классов, так как побочные эффекты этих действий не влияют на конфликты между правилами. Однако для некоторых сценариев имеет смысл сохранить различение между действиями при анализе множеств, которые относятся к конкретному действию.

В работе Diekmann и др.~\cite{diekmann2018verified} было показано, что при формализации правил iptables достаточно только двух действий --- \emph{allow} и \emph{deny}, так как действия с побочным эффектом (log, count и др.) не влияют на судьбу пакета, а цепочки всегда можно развернуть в отсутствие петель; таким образом, исходный набор можно привести к набору правил с только двумя действиями --- \emph{allow} и \emph{deny}. Это корректно с точки зрения анализа правил фильтрации, так как системе анализа не интересен побочный эффект правила.

\subsection{Качество конфигураций в реальной эксплуатации}
\label{sec:Section2:empirical}

Прежде чем переходить к определению видов различных ошибок, стоит установить природу их возникновения. Можно представить, что однажды аккуратно спроектированный набор правил эксплуатируется без накопления ошибок. Однако эмпирические данные, собранные за последние два десятилетия, количественно фиксируют закономерности деградации конфигураций.

Одной из ранних работ в данной области остаётся исследование Wool~\cite{wool2004quantitative}. В ней автор собрал 37 наборов правил межсетевых экранов Check Point FireWall-1 из крупных организаций разных секторов и оценил их по 12 типам ошибок, описанных в работе. Главный количественный вывод: более 90\,\% обследованных межсетевых экранов содержали хотя бы одну грубую ошибку, причём девять из двенадцати ошибок встречались более чем в половине конфигураций.
Помимо собранной статистики Wool ввёл метрику сложности конфигурации $RC$:

\begin{equation*}
    RC = N_{\text{rules}} + N_{\text{objects}} + \frac{N_{\text{interfaces}}(N_{\text{interfaces}} - 1)}{2}
\end{equation*}

где $N_{\text{rules}}$, $N_{\text{objects}}$, $N_{\text{interfaces}}$ --- количество правил, объектов и интерфейсов, используемых в наборе соответственно. Он показал, что число ошибок зависит от сложности конфигурации логарифмически: $\#\text{errors} \approx \ln(RC) + 1{,}5$.

Через шесть лет Wool~\cite{wool2010trends} повторил исследование на расширенной выборке из 84 наборов правил Check Point и Cisco PIX, разделил ошибки на 36 категорий по четырём группам (входящий, исходящий, внутренний трафик, заведомо рискованные правила) и подтвердил корреляцию числа ошибок и сложности конфигурации. В этот раз уже потребовалась вендоро-независимая метрика сложности конфигурации межсетевого экрана $FC$:

\begin{equation*}
FC_p = N_{\text{lines}} - 50, \qquad FC_c = N_{\text{rules}} \cdot N_{\text{interfaces}} + N_{\text{objects}},
\end{equation*}

где $N_{\text{lines}}$ --- число строк конфигурации правил Cisco PIX, $N_{\text{rules}}$, $N_{\text{interfaces}}$, $N_{\text{objects}}$ --- количество правил, интерфейсов и объектов соответственно для Check Point.

Главным результатом повторного исследования стал факт --- межсетевые экраны более новых версий не позволяют администраторам совершать меньшее число ошибок. Таким образом, от необходимости в аудите конфигураций невозможно избавиться переходом на новую платформу.

Аналогичные результаты фиксируются и в индустриальных источниках: \emph{Verizon Data Breach Investigations Report}~\cite{verizon2010data} систематически фиксирует ошибки конфигурации сетевых средств защиты как одну из устойчивых причин утечек данных наряду с эксплуатацией уязвимостей и социальной инженерией. Таким образом, задача автоматизированного аудита правил мотивируется не только академическими метриками, но и реальными издержками эксплуатации.

В общем итоге имеем три важных вывода о качестве политик безопасности и правил фильтрации в частности:

\begin{enumerate}
    \item Ошибки в наборах правил накапливаются гарантированно: их доля в обследованных конфигурациях достаточно велика, чтобы рассматривать ручную проверку как недостаточную.
    \item Основным предиктором числа ошибок является структурная сложность конфигурации, а не только число правил само по себе.
    \item Проблема наличия ошибок в наборах не решается сменой платформы.
\end{enumerate}

\subsection{Типы ошибок и аномалий в наборах правил}
\label{sec:Section2:anomalies}

Понятие ошибок и аномалий занимает ключевое место в системах анализа правил. Разные авторы вводили близкие по смыслу, но не всегда совпадающие определения, поэтому на данный момент существует несколько различных классификаций. За основу взята модель, представленная в работе Al-Shaer и Hamed~\cite{al2003firewall} и формализована Yuan и др. в работе~\cite{yuan2006fireman}. Также Al-Shaer и Hamed расширили свою таксономии за пределами правил фильтрации \cite{hamed2006taxonomy}: расширили на политики IPSec и QoS.

\subsubsection*{Отношения между парами правил}

Al-Shaer и Hamed~\cite{al2003firewall} вводят формальные отношения между парами правил ($R_x$ и $R_y$) с помощью полей ($R_x[i]$ и $R_y[i]$) из five-tuple:
\begin{itemize}
    \item \textbf{Полностью различимые (Completely disjoint, CD)} --- ни одно поле $R_x$ и $R_y$ не совпадает
    \item \textbf{Точно совпадающие (Exactly matching, EM)} --- все поля совпадают; такая пара представляет вырожденный случай, при котором правило с большим индексом гарантированно либо избыточно (если действия совпадают), либо затенено (если различаются).
    \item \textbf{Полностью включает (Inclusively matching, IM)} --- все поля одного правила являются подмножествами этих полей другого правила, но при этом хотя бы одно поле не совпадает.
    \item \textbf{Частично соответсвует (Partially matching, PM)} --- если хотя бы над одним полем выполенено отношения строгого подмножества, надмножеством или равенства, и существует другое поле, над которым это не выполнено.
    \item \textbf{Коррелирующие (Correlated, C)} --- по некоторым поля правило является строгим подмножеством или равенство, а остальным поля являются строгим надмножеством.
\end{itemize}

Авторы доказывают теорему о полноте этой таксономии: других отношений между парами 5-tuple правил не существует. Та же таксономия (IM/EM/PM/CD/C) воспроизводится в более поздних обзорах, в частности у Ramesh и др.~\cite{ramesh2024exploring}.

Радомский~\cite{radomskiy2011security} расширяет приведённую таксономию добавлением отношения \textbf{соседства (adjacency)}: два правила соседствуют, если существует единственное поле, которое по которому можно объединить в связный интервал, остальные поля совпадают. Это отношение не входит в полные пять классов Al-Shaer и принципиально необходимо для определения union-аномалии (см. ниже): без явного отношения соседства два смежных правила выглядят как пара частично совпадающих, над которыми не всегда можно определить операцию объединения в одно правило.

\subsubsection*{Классические аномалии}

На основе перечисленных отношений Al-Shaer и Hamed определяют четыре класса аномалий между правилами $R_x$ и $R_y$ при $x < y$ (то есть $R_x$ предшествует $R_y$ в списке); используется бинарная система действий:
\begin{itemize}
    \item \textbf{Затененние (Shadowing)} --- $R_x$ полностью включает или совпадает с $R_y$; $R_x$ и $R_y$ имеют разные действия; $R_y$ становится недостижимым: ни один пакет не дойдёт до $R_y$, поскольку будет обработан $R_x$. Это явная ошибка администратора.
    \item \textbf{Избыточность (Redundancy)} --- $R_x$ полностью включает или совпадает с $R_y$; $R_x$ и $R_y$ имеют одинаковые действия; удаление $R_y$ не изменит политику с точки зрения множества прошедших и заблокированных пакетов. Это не является ошибкой в строгом смысле, но указывает на избыточность набора. Acharya и др.~\cite{acharya2006traffic} выделяют два вида избыточности: внешнюю (\emph{external redundancy}), в смысле выше, и внутренюю (\emph{internal redundancy}) в тех случаях, когда в поле можно указать некоторый набор значений, и он избыточен уже, или написан неоптимально с точки зрения объединения. Внутренняя избыточность также не нарушает семантику, но затрудняет ручной аудит.
    \item \textbf{Обобщение (Generalization)} --- $R_x$ является полным включением $R_y$ с противоположным действием. В общем случае не ошибка: нередко это сознательный паттерн дизайна <<разрешить общее, кроме конкретного>> (например, разрешить весь HTTP, кроме ряда подозрительных хостов). Поэтому инструмент должен сообщать о таких парах как о подозрительных, но не помечать как ошибку.
    \item \textbf{Корреляция (Correlation)} --- $R_x$ и $R_y$ коррелирует, т.е. они имееют пересечение по некоторому множестве пакетов, но ни одно не является подмножеством другого;имеют разные действие. Результат обработки таких пакетов зависит от порядка правил, что почти всегда является непреднамеренным, и такие правила лучше переписать.
\end{itemize}

\subsubsection*{Формализация в терминах множеств}

Yuan и др.~\cite{yuan2006fireman} в инструменте FIREMAN переформулируют те же аномалии в терминах теоретико-множественной семантики набора правил. В работе рассматриваются не только обычные линейные списки правил, но и механизм цепочек, т.е. система действий с переходами (jump, return). Анализатор не работает с цепочками напрямую, происходит операция выделения всех возможных путей обхода связанных списков правил. Для построения путей добавляются специальные правила с действием pass, которое означает случай, когда при обходе мы остались на этом путе.

Так, для каждой позиции $j$ для каждого пути обхода поддерживается состояние $\langle A_j, D_j, F_j\rangle$, где $A_j$ --- множество пакетов, уже принятых предшествующими правилами, $D_j$ --- множество отброшенных, $F_j$ --- множество пакетов, ушедших на альтернативный путь; множество пакетов, доступных для рассмотрения текущим правилом, определяется как $R_j = I \setminus (A_j \cup D_j \cup F_j)$, где $I$ --- общее множество входящих пакетов.

Тогда для правила $\langle P_j, \text{accept}\rangle$ различают следующие случаи:
\begin{itemize}
    \item $P_j \subseteq R_j$ --- правило корректное; оно определяет действие для нового множества пакетов, не пересекающегося с предшествующими правилами;
    \item $P_j \cap R_j = \emptyset$ --- \emph{masked rule}: правило не сработает ни на одном пакете (ошибка). Подслучаи:
    \begin{itemize}
        \item $P_j \subseteq D_j$ --- \emph{shadowing}: правило должно принимать пакеты, которые уже отброшены предыдущими правилами;
        \item $P_j \cap D_j = \emptyset$ --- \emph{redundancy}: все пакеты $P_j$ уже приняты предшествующими правилами либо ушли на альтернативный путь;
        \item иначе --- одновременно \emph{redundancy} и \emph{correlation}: часть пакетов уже отброшена, остальные приняты или ушли на альтернативный путь;
    \end{itemize}
    \item $P_j \not\subseteq R_j \wedge P_j \cap R_j \neq \emptyset$ --- \emph{partially masked rule}: правило сработает не на всех заявленных пакетах. Возможны (не взаимоисключающие) подслучаи:
    \begin{itemize}
        \item $P_j \cap D_j \neq \emptyset$ --- \emph{correlation} (предупреждение): часть пакетов, которые правило намерено принять, уже отброшена предшествующими правилами;
        \item существует более раннее правило $\langle P_x, \text{deny}\rangle$ при $P_x \subseteq P_j$ --- \emph{generalization} (предупреждение): правило $j$ обобщает правило $x$ с противоположным действием;
        \item $P_j \cap A_j \neq \emptyset$ и существует более раннее $\langle P_x, \text{accept}\rangle$ при $P_x \subseteq P_j$ --- \emph{redundancy} (ошибка) правила $x$: его можно удалить, так как покрываемые им пакеты всё равно будут приняты текущим правилом.
    \end{itemize}
\end{itemize}
Симметричная таблица справедлива для правил с действием \emph{deny} (заменой $A_j \leftrightarrow D_j$).

Эта формулировка эквивалентна определениям Al-Shaer для пар правил, но дополнительно охватывает случаи, когда аномалия формируется \emph{совместным} действием нескольких правил, --- именно за счёт того, что состояния $A_j$, $D_j$ агрегируют эффект всех предшествующих правил, а не одного.

Помимо аномалий в рамках одного набора правил, FIREMAN формализует классы нарушений уровня сети целиком: \emph{policy violation} (расхождение фактической политики с заданным белыми и черными списками), \emph{inefficiency} (избыточность и неоптимальность набора правил без нарушения семантики), а для распределённых конфигураций --- \emph{inter-firewall inconsistency} (последовательные межсетевые экраны на одном маршруте принимают противоречащие решения для одного и того же пакета) и \emph{cross-path inconsistency} (разные маршруты между одной парой узлов приводят к разным итоговым решениям из-за различающихся множеств фильтров на пути). В настоящей работе эти классы не рассматриваются: корректный анализ достижимости в сети требует учёта маршрутизации, NAT, VPN-туннелей и динамики управляющего плана, то есть факторов вне правил фильтрации, и относится к отдельному направлению верификации сети --- Header Space Analysis~\cite{kazemian2012header}, APKeep~\cite{zhang2020apkeep}, Katra~\cite{beckett2022katra} и других работ по анализу достижимости.

\subsubsection*{Структурные свойства политики: consistency и completeness}

Параллельно с попарной классификацией Al-Shaer, в работах Gouda и Liu~\cite{gouda2004firewall, gouda2007structured} формулируются два структурных свойства политики, нарушения которых сами по себе образуют классы ошибок. Свойства определяются для FDD, но содержательно характеризуют любую политику фильтрации:
\begin{itemize}
    \item \textbf{Consistency} --- для каждого пакета определено единственное решение: $f.\text{accept} \cap f.\text{discard} = \emptyset$. Нарушение consistency означает, что один и тот же пакет в одних правилах разрешён, а в других запрещён; в спискообразных рулсетах это противоречие маскируется first-match-семантикой, и проявляется как зависимость поведения от порядка правил --- то есть совпадает с классической correlation- или shadowing-аномалией. На уровне политики целиком consistency-нарушение означает невозможность дать однозначный ответ <<разрешён ли пакет $p$>> без знания порядка обработки.
    \item \textbf{Completeness} --- любой возможный пакет покрыт каким-либо правилом: $f.\text{accept} \cup f.\text{discard} = \Sigma$. Нарушение completeness означает наличие <<серой зоны>> --- пакетов, не сопоставленных ни одному пользовательскому правилу. На реальных межсетевых экранах эта зона неявно поглощается default-правилом (\emph{deny all} или \emph{accept all}), поэтому incompleteness не приводит к runtime-сбою, но скрывает за умолчанием намерения администратора и затрудняет аудит.
\end{itemize}
Любая политика, описанная в виде явного FDD, по построению удовлетворяет обоим свойствам; для произвольного спискового рулсета они должны быть проверены анализатором как самостоятельный класс структурных проверок.

\subsubsection*{Расширения: union, irrelevance, групповые аномалии}

Ограничение попарной классификации в том, что аномалия может быть результатом совместного действия трёх и более правил. В работе Chomsiri и Pornavalai~\cite{chomsiri2006firewall} (Raining 2D-Box Model) эта проблема явно поставлена и решена для случая, когда shadowing формируется коллективно. Радомский~\cite{radomskiy2011security} систематизирует подобные ситуации и дополнительно вводит:
\begin{itemize}
    \item \textbf{Union anomaly} --- два или более adjacent-правил с одинаковым решением, которые можно объединить в одно без изменения семантики. Возможность такого объединения опирается на отношение adjacency, введённое выше; без него union-аномалия не определима.
    \item \textbf{Irrelevance anomaly} --- правило, которому не сопоставляется ни один достижимый пакет (например, из-за маршрутизации или предшествующего перехвата трафика); в распределённых межсетевых экранах эта ситуация впервые описана Mohamed и др. (по~\cite{radomskiy2011security}).
\end{itemize}
Радомский также предлагает методически важное разделение всех аномалий на две категории: \emph{errors} --- аномалии, для которых возможно автоматическое исправление без потери семантики (shadowing, redundancy, internal redundancy), и \emph{warnings} --- аномалии, требующие подтверждения администратора, поскольку могут быть осознанным паттерном дизайна (generalization, correlation, union). Это разделение принимается в настоящей работе как принцип формирования отчёта анализатора.

Следует отметить, что union-аномалии в современной практике эксплуатации сетей представляют меньший интерес, чем кажется на первый взгляд. Организация zero-trust-инфраструктуры предполагает построение политик от объектов и привязку правил к заявкам пользователей; объединение таких правил без потери семантики невозможно, поскольку каждое из них несёт метаинформацию (номер заявки, ответственное лицо, срок действия), которую агрегированное правило сохранить не может. Поэтому в настоящей работе union-анализ рассматривается как опциональная возможность, выводимая в виде предупреждения, а не как обязательное преобразование.

\subsubsection*{Аномалии, выявляемые только по трафику}

Особое место занимают классы аномалий, не выявляемые статическим анализом и обнаружимые только при сопоставлении политики с реально наблюдаемым трафиком. Golnabi и др.~\cite{golnabi2006analysis} вводят два таких класса: \emph{блокировка существующего сервиса} (политика блокирует трафик, который реально требуется и наблюдается в журналах) и \emph{разрешение трафика к несуществующему сервису} (правило разрешает трафик, для которого не существует целевой системы). Помимо этого, авторы вводят понятия \textbf{decaying rules} --- правил, чей трафик статистически падает со временем, что указывает на потерю их актуальности, и \textbf{dominant rules} --- правил, обрабатывающих основную долю трафика и потому критически важных для производительности. К той же категории по своей природе относятся \textbf{stale rules} --- правила, оставшиеся от прежних конфигураций сервисов, давно выведенных из эксплуатации; такая аномалия зафиксирована Diekmann и др.~\cite{diekmann2018verified} в одном из реальных кейсов (правило для удалённого SSH-доступа, оставшееся в политике после окончания соответствующего эксперимента). Все эти классы требуют данных журналов или внешнего инвентаря сервисов; соответствующие методы рассматриваются ниже в~\ref{sec:Section2:logs}.

\subsubsection*{Аномалии за пределами одного межсетевого экрана}

Часть классов аномалий относится не к политике одного устройства, а к сети в целом: \textbf{forwarding loops} (циклы передачи), \textbf{reachability failures} (потеря связности между легитимными узлами), \textbf{traffic isolation/leakage} (нарушение изоляции slice'ов или подсетей). Эти классы формализованы в Header Space Analysis Kazemian и др.~\cite{kazemian2012header} и в работах семейства VeriFlow и относятся к задаче верификации связности сети, а не аудита одного межсетевого экрана. В настоящей работе они находятся за пределами основного scope; высокоуровневые требования к связности при необходимости могут быть проверены поверх той же модели правил, как обсуждается в~\ref{sec:Section2:distributed}.

\subsection{Методы анализа правил фильтрации}
\label{sec:Section2:methods}

\subsubsection{Попарный анализ и геометрические интерпретации}
\label{sec:Section2:pairwise}

Исторически первое формальное направление --- попарный анализ правил с целью обнаружения аномалий. Помимо введения упомянутой выше таксономии, Al-Shaer и Hamed~\cite{al2003firewall} реализовали инструмент Firewall Policy Advisor, который вместо наивного попарного перебора $O(n^2)$ строит многоуровневое \emph{policy tree} (protocol $\to$ src\_ip $\to$ src\_port $\to$ dst\_ip $\to$ dst\_port). Каждое правило соответствует пути от корня к листу; конфликты выявляются как множественные пути, проходящие через одни и те же узлы с разными действиями. Помимо диагностики, инструмент поддерживает \emph{anomaly-free editing}: при добавлении нового правила алгоритм автоматически вычисляет позицию вставки, не порождающую новых конфликтов.

Близкие по замыслу работы Hari, Suri и Parulkar~\cite{hari2000detecting} и Eppstein и Muthukrishnan~\cite{eppstein2000internet} опираются на геометрическую интерпретацию правила как $d$-мерного прямоугольника в пространстве заголовков пакетов и сводят задачу обнаружения конфликтов к задачам вычислительной геометрии. Hari и др. формально характеризуют условия конфликта в фильтрах с произвольным числом полей и доказывают, что схемы разрешения конфликтов на основе явного приоритета или порядка фильтров \emph{принципиально неполны}: существуют наборы правил, для которых никакая перестановка не даёт желаемого поведения. В качестве решения они предлагают идею \emph{resolve-фильтров} --- добавление дополнительного правила, явно покрывающего область пересечения. Eppstein и Muthukrishnan, в свою очередь, дают существенное теоретическое улучшение: для $d=2$ обнаружение всех конфликтов выполняется за $O(n^{3/2})$ при линейной памяти --- против тривиального $O(n^2 \log n)$, --- а классификация пакетов выполняется за $O(\log\log n)$ при $O(n^{1+o(1)})$ памяти. Эти результаты задают теоретический предел для производительности pairwise-анализа, но обе работы ограничиваются обнаружением единичных конфликтов между парами правил и не дают целостной картины политики.

Принципиальное ограничение попарного подхода --- неспособность обнаруживать аномалии, формируемые совместным действием трёх и более правил. Это ограничение и привело к появлению символьных методов, рассматривающих набор правил как единое целое.

\subsubsection{Символьные методы анализа политик}
\label{sec:Section2:symbolic}

Параллельно с pairwise-анализом сформировалось семейство \emph{query-based} инструментов: Fang Mayer, Wool и Ziskind~\cite{mayer2000fang} (IEEE S\&P 2000), а также позднее коммерческий Lumeta. Fang парсит реальные конфигурационные файлы межсетевых экранов разных вендоров (Cisco, Lucent), строит общую внутреннюю модель топологии и rule-base'ов и позволяет администратору задавать запросы вида <<разрешает ли политика сервис $s$ от $a$ к $b$?>>; учитываются маршрутизация, NAT, spoofing-сценарии. Эта работа впервые сформулировала ключевую терминологию (\emph{gateway, interface, zone, host-group, service}), оставшуюся стандартной в области, и первой реализовала passive-аналитический инструмент, не требующий активного зондирования сети. Однако Fang проверяет лишь конкретные запросы пользователя, а не всё пространство пакетов: пропущенные сценарии остаются без покрытия.

Революционной для области стала работа Yuan и др.~\cite{yuan2006fireman} \emph{FIREMAN: A Toolkit for FIREwall Modeling and ANalysis} (IEEE S\&P 2006). Авторы предложили рассматривать набор правил как программу, которую можно проанализировать методами символьного исполнения и model checking, получив полную картину разрешаемого и блокируемого трафика без явного перебора пакетов. Каждое правило записывается в виде $\langle P, \text{action}\rangle$, где $P$ --- предикат над полями пакета, а действие принадлежит множеству $\{\text{accept},\text{deny},\text{chain }Y,\text{return}\}$. Правила обновления состояния $\langle A,D,F\rangle$, введённого выше:
\begin{align*}
\langle A,D,F\rangle, \langle P,\text{accept}\rangle &\vdash \langle A \cup (R \cap P), D, F\rangle, \\
\langle A,D,F\rangle, \langle P,\text{deny}\rangle &\vdash \langle A, D \cup (R \cap P), F\rangle, \\
\langle A,D,F\rangle, \langle P,\text{pass}\rangle &\vdash \langle A, D, F \cup (R \cap \neg P)\rangle.
\end{align*}
Финальные множества $A_{ACL} = \bigcup_i A_{path_i}$ и $D_{ACL} = \bigcup_i D_{path_i}$ удовлетворяют свойству $A_{ACL} \cup D_{ACL} = I_{ACL}$. Все проверки на аномалии (приведённые в~\ref{sec:Section2:anomalies}) сводятся к булевым операциям над предикатами. Ключевой инженерной находкой FIREMAN стало использование двоичных решающих диаграмм (BDD) из библиотеки BuDDy для канонического представления предикатов; операции пересечения, объединения и проверки включения сводятся к стандартным операциям над BDD. Поскольку число состояний межсетевого экрана конечно, символьное model checking гарантирует одновременно полноту и корректность.

FIREMAN рассматривает не только одиночный межсетевой экран, но и распределённую сеть. Для последовательно соединённых межсетевых экранов справедливы соотношения $A_{\text{serial}} = \bigcap A_{acl}$, $D_{\text{serial}} = \bigcup D_{acl}$, а для параллельно соединённых --- $A_{\text{parallel}} = \bigcup A_{acl}$, $D_{\text{parallel}} = \bigcap D_{acl}$. Эти выражения позволяют построить ACL-дерево с корнем в защищаемой сети и сравнить агрегированную политику с заданными black-list и white-list для обнаружения нарушений политики и cross-path-аномалий. На реальных промышленных конфигурациях (PIX1 на 249 правил, PIX2 на 36, BSD1 на 94) FIREMAN обнаружил суммарно несколько десятков ранее неизвестных нарушений; производительность алгоритма для intra-firewall-анализа --- $O(n)$ при использовании BDD против $O(n^2)$ у Al-Shaer; набор из 800 правил обрабатывается за три секунды.

Ограничения FIREMAN, существенные для настоящей работы: анализ выполняется без учёта состояний (stateless); whitelist должен задаваться вручную; в распределённой схеме предполагается, что любой топологически возможный путь может быть выбран маршрутизацией, что в реальных сетях даёт ложные срабатывания; категории correlation и generalization выводятся как warning; инструмент диагностирует, но не исправляет ошибки; расширение модели на современные NGFW и нестандартные match-условия не предусмотрено.

Альтернативное символьное представление политики было предложено в более ранних работах Gouda и Liu~\cite{gouda2004firewall} (\emph{Firewall Decision Diagram}, ICDCS 2004). FDD --- это направленный ацикличный граф над полями $F_0, \dots, F_{n-1}$, удовлетворяющий условиям: существует единственная корневая вершина; внутренние вершины помечены полями, листья --- решениями; ни одно поле не повторяется на одном пути от корня к листу; каждое ребро помечено непустым подмножеством области соответствующего поля; для каждой вершины множества меток исходящих рёбер не пересекаются (consistency) и в объединении дают всю область поля (completeness). Каждый путь от корня к листу определяет одно правило вида $F_0 \in S_0 \wedge \dots \wedge F_{n-1} \in S_{n-1} \to \text{decision}$. Из определения сразу следует ключевая теорема: $f.\text{accept} \cap f.\text{discard} = \emptyset$ и $f.\text{accept} \cup f.\text{discard} = \Sigma$, то есть FDD по построению не допускает ошибок consistency и completeness.

FDD-ориентированный подход выглядит привлекательно, но имеет несколько ограничений, существенных для настоящей работы. Во-первых, диаграмма строится на этапе проектирования межсетевого экрана с нуля; адаптация существующих правил с произвольными match-условиями к FDD требует значительной предобработки и не всегда возможна. Во-вторых, представление каждого поля сводится к интервальному (или префиксному) множеству целых, что несовместимо с современным разнообразием match-условий: списки приложений, URL-категории, временные интервалы, доменные имена. В-третьих, эффективное представление достижимости (reachability) в FDD получается громоздким, поскольку требует явного обхода всех путей. Поэтому FDD в настоящей работе используется как теоретический ориентир, но не как основная структура данных.

В работе Liu~\cite{liu2008formal} (ICC 2008) на основе FDD построен формально верифицируемый алгоритм проверки заданных свойств политики. Свойство задаётся в виде property rule $\bigwedge_i F_i \in Y_i \to \text{decision}$ и проверяется обходом FDD; на реальных межсетевых экранах в 87 и 661 правил верификация занимает миллисекунды. Acharya и Gouda~\cite{acharya2011firewall} (INFOCOM 2011) формализуют фундаментальный результат: задачи верификации заданного свойства и обнаружения избыточности правила эквивалентны. Авторы строят алгоритмы $V \to R$ и $R \to V$ с сохранением асимптотической сложности. Этот результат является важным теоретическим основанием: достаточно построить эффективный движок для одной из этих задач, чтобы автоматически решать обе.

\subsubsection{Логические методы и SMT-подходы}
\label{sec:Section2:smt}

Параллельно с символьными методами развивались подходы, опирающиеся на SAT- и SMT-решатели. Работа Zhang, Mahmoud, Malik и Narain~\cite{zhang2012verification} (ICNP 2012) формулирует задачу верификации эквивалентности и включения межсетевых экранов как задачу выполнимости булевых формул и решает её через MiniSat. Каждое правило кодируется тернарным массивом битов; для битового положения $j$ правила $i$ с метаусловием $r_{i,j}$ вводится формула
\begin{equation*}
k_{i,j} = \begin{cases} b_j, & r_{i,j} = 1 \\ \neg b_j, & r_{i,j} = 0 \\ \mathrm{TRUE}, & r_{i,j} = x \end{cases}.
\end{equation*}
Match-переменная $m_i$ выражается через конъюнкцию $k_{i,j}$ и условие <<ни одно из предыдущих правил не сработало>>, а выход политики --- через дизъюнкцию match-переменных permit-правил. Эквивалентность двух политик $F_1$ и $F_2$ проверяется через выполнимость формулы $F_1(B) \oplus F_2(B)$; включение --- через $\neg(F_1(B) \to F_2(B))$. Обнаружение избыточности производится введением управляющего бита $o_i$ для каждого правила: при $o_i = 1$ правило отключено, и алгоритм жадно ищет максимальное множество правил, которые можно отключить без изменения семантики.

Авторы также предложили формализацию задачи синтеза минимального межсетевого экрана как QBF-задачи $\exists V \forall B : \neg(F(B) \oplus F^s(B,n,V))$, однако эксперимент показал, что даже простейшие случаи ($N_{\text{input}} = 25$) выходят за пределы возможностей современных QBF-решателей. SAT-методы для верификации, напротив, успешно работают на конфигурациях до 26 000 правил со временем порядка минут.

Близкий по духу инструмент --- Margrave Nelson, Barratt, Dougherty, Fisler и Krishnamurthi~\cite{nelson2010margrave} (LISA 2010) --- использует формализм первого порядка и SAT-решатель Kodkod. Margrave поддерживает запросы на разных уровнях абстракции (правило, фильтр, межсетевой экран, сеть межсетевых экранов), а его выход --- сценарии (concrete witnesses), показывающие конкретные пакеты, удовлетворяющие условию запроса. Margrave принципиально применим только к Cisco IOS, поскольку компилятор для других систем не реализован. Также важной идеей этой работы является поддержка change-impact-анализа в виде первоклассного запроса: $(P_1.\text{Permit} \wedge \neg P_2.\text{Permit}) \vee (\neg P_1.\text{Permit} \wedge P_2.\text{Permit})$.

Алгоритмы change-impact, занимающие центральное место в работах Liu~\cite{liu2008firewall} (ACM TOIT 2012) и Margrave, в настоящей работе \emph{не входят в основной функционал}. Причина в том, что change-impact-анализ относится к категории pre-change анализа: администратор перед внесением изменения получает картину последствий. Разрабатываемая система ориентирована на post-change аудит: на регулярную проверку уже сложившейся, эксплуатируемой конфигурации с целью обнаружения накопленных ошибок. Технически change-impact алгоритм может быть включён как опциональное расширение, но смысловая ось разрабатываемого инструмента сосредоточена на статическом обнаружении аномалий в существующей политике.

Из современных работ этой линии следует отметить Togay и др.~\cite{togay2021firewall} (IEEE Transactions on Reliability 2021), строящих фреймворк FPAD на Prolog для непрерывного мониторинга правил. Эту работу можно рассматривать как развитие линии Al-Shaer/FIREMAN с применением современных формальных методов и подходов машинного обучения.

\subsubsection{Формально верифицированный анализ реальных конфигураций}
\label{sec:Section2:verified}

К наиболее значимым работам последнего десятилетия относится исследование Diekmann и др.~\cite{diekmann2018verified} \emph{Verified iptables Firewall Analysis and Verification} (Journal of Automated Reasoning, 2018). Эта работа объединяет несколько важных идей и в значительной мере определила архитектуру разрабатываемого в настоящей ВКР средства анализа.

Авторы исходят из того, что между сложностью реальной iptables-конфигурации и упрощённой моделью академических анализаторов существует значительный разрыв. Iptables в современных версиях поддерживает свыше 60 типов match-условий и более 200 опций; user-defined chains, conntrack-состояния, ratelimit-конструкции, строковые матчеры, conntrack-helpers --- ни один из публично доступных инструментов не обрабатывает их в полном объёме. Авторы экспериментально продемонстрировали, что инструмент ITVal на простом наборе правил NAS-устройства Synology выдаёт неверный результат, а другие инструменты вообще завершаются ошибкой.

Решение, предложенное Diekmann и др., состоит в том, чтобы построить формально верифицированный <<препроцессор>>: серию преобразований над набором правил iptables, каждое из которых сохраняет (или контролируемо приближает) семантику и сводит сложный набор к простой list-модели, доступной для анализа. Все преобразования и леммы реализованы в Isabelle/HOL и формально доказаны. Архитектура работы включает:
\begin{enumerate}
    \item \textbf{Формальную семантику iptables} в виде индуктивного предиката $\Gamma, \gamma, p \vdash \langle rs, t\rangle \Rightarrow t'$, где $\Gamma$ --- окружение пользовательских цепочек, $\gamma$ --- функция матчинга, $p$ --- пакет, $rs$ --- список правил, $t$ и $t'$ --- состояния (Decision: Accept/Drop/Undecided).
    \item \textbf{Chain unfolding} --- развёртывание call/return пользовательских цепочек в плоский список, формально доказывающее теорему о семантической эквивалентности.
    \item \textbf{Тернарную логику для неизвестных match-условий}: если матчер $m$ не распознан, он возвращает значение $\bot$. Поведение для пакетов, на которых матчер не определён, аппроксимируется двумя замыканиями --- upper closure ($\bot \to \top$) и lower closure ($\bot \to \bot$). Тогда множество пакетов, гарантированно блокируемых реальной политикой, не превышает множества, разрешённого upper closure, а множество, гарантированно разрешённых, не превышает разрешённого lower closure.
    \item \textbf{Алгебру match-выражений} --- нормализация в DNF, удаление тавтологий и противоречий, сведение конъюнкций IP-диапазонов к канонической форме, замена interface-условий IP-условиями через таблицу маршрутизации.
    \item \textbf{Простую модель межсетевого экрана} --- список правил вида $\langle in\_iface, out\_iface, src\_ip, dst\_ip, proto, src\_port, dst\_port, action\rangle$ без call/return и без неизвестных матчеров; эта модель служит входом для анализаторов.
    \item \textbf{Аналитические методы}: разбиение пространства IP-адресов на классы эквивалентности по политике и построение сервисных матриц.
    \item \textbf{Самостоятельный инструмент fffuu}, экспортированный из Isabelle в Standard ML и работающий с выводом \texttt{iptables-save}.
\end{enumerate}

В рамках экспериментальной части работы авторы собрали крупнейший в академии набор реальных iptables-конфигураций. На NAS-устройстве Synology, на котором ITVal выдаёт некорректный результат (объединяет всё пространство IP в один класс), fffuu правильно выделяет локальную сеть. На firewall H 2003 года upper closure обнаружил, что фактически принимается весь трафик: при множестве permit-правил итоговая политика по умолчанию --- accept, что является серьёзной ошибкой; lower closure показал отсутствие защиты от MAC-spoofing-атак.

Из работы Diekmann и др. в настоящем исследовании заимствуются следующие принципиальные идеи. Во-первых, разделение анализа на две фазы --- нормализация (упрощение и приведение к канонической форме) и собственно анализ --- позволяет отделить вендоро-зависимую часть от анализа. Во-вторых, тернарная логика для неизвестных match-условий с гарантированной upper/lower closure --- крайне привлекательный подход, поскольку позволяет работать со \emph{всеми} реальными конфигурациями, даже если часть match-условий не распознана инструментом. В-третьих, сервисные матрицы --- наглядный формат отчёта пользователю.

Принципиальное ограничение Diekmann --- инструмент привязан к iptables. Хотя авторы в обсуждении указывают, что методология применима к Cisco IOS ACL и OpenFlow, формальная семантика и Isabelle/HOL-формализация для других систем заново не построены. Кроме того, авторы рассматривают только filter-таблицу, не покрывая NAT и mangle, а также не моделируют динамику соединений. Формально верифицированный подход Diekmann показывает, что аппроксимация иногда теряет информацию: для firewall H upper и lower closures дают <<всё>> и <<ничего>>. Для разрабатываемого инструмента это указывает на необходимость рассматривать вместо полностью неизвестных match-условий хотя бы пятёрку (5-tuple) полей правила, остальные поля считая допускающими любое значение, --- так результат аппроксимации остаётся информативным.

\subsubsection{Верификация распределённых межсетевых экранов и SDN}
\label{sec:Section2:distributed}

Отдельной ветвью являются работы по верификации сетей межсетевых экранов и программно-конфигурируемых сетей. Работа Kapus~\cite{kapus2023improved} (IEEE Access 2023) и её предшественник 2020 года используют формализм TLA$^+$ и model checker TLC для проверки соответствия распределённой SDN-инфраструктуры заданной монолитной политике. В улучшенной версии 2023 года авторы устраняют переменные состояния коммутаторов, моделируют коммутатор как функцию <<пакет, входной порт $\to$ множество выходных портов>> и используют рекурсивную обходную процедуру по DAG. Это позволяет сократить время верификации с часов до секунд на тех же тестовых сетях.

Аналогичные работы по reachability-анализу --- Header Space Analysis Kazemian, Varghese и McKeown~\cite{kazemian2012header} (NSDI 2012), VeriFlow и др. --- образуют семейство методов, ориентированных на верификацию связности сетевых политик в крупных дата-центрах. Header Space Analysis примечательна протокол-агностичным подходом: заголовок пакета рассматривается как точка в пространстве $\{0, 1\}^L$ без априорного знания структуры полей, а каждое сетевое устройство моделируется \emph{transfer function} $T : P \to 2^P$, преобразующей подмножества заголовков. Алгебра пересечений и объединений wildcard-выражений позволяет статически проверять reachability, обнаруживать forwarding-loop'ы и нарушения изоляции slice'ов. На реальной кампусной сети Стэнфорда все forwarding-loop'ы найдены за время порядка 10 минут, проверка достижимости между двумя подсетями --- за 13 секунд.

Параллельно reachability-методам развивалась линия визуального аудита: Mansmann, Göbel и Cheswick~\cite{mansmann2012visual} (VizSec 2012) демонстрируют, что наборы правил масштаба нескольких сотен и тысяч строк не поддаются ручному пересмотру и предлагают визуальные представления, помечающие зоны высокой плотности конфликтов. С точки зрения настоящей работы это указывает на отдельное требование к выходу анализатора --- не только список аномалий, но и форма представления, удобная для оператора, обозревающего политику в целом.

Существенный вклад в эту линию вносит работа Bringhenti и др.~\cite{bringhenti2020improving} (IEEE TNSM 2021), где предлагается моделировать функции не на уровне отдельных пакетов, а на уровне \emph{packet classes} --- максимальных подмножеств, на которых все network functions ведут себя одинаково. Это позволяет существенно ускорить SMT-верификацию reachability в виртуализированных сетях с stateful-функциями (NAT, conntrack-firewall) и поддержать сложные service function chains.

В целом эти работы технически сильны, но решают задачу, которая для разрабатываемого инструмента не является приоритетной. Современная эксплуатация корпоративных сетей всё чаще тяготеет к консолидации защиты на нескольких узловых точках (центральный NGFW, периметровые маршрутизаторы) вместо распределённых межсетевых экранов. Поэтому в настоящей работе анализ распределённых межсетевых экранов рассматривается как возможное расширение, тогда как высокоуровневые требования к связности (<<какие подсети могут связываться с какими>>) могут быть проверены методами анализа достижимости поверх той же модели правил, которая используется для поиска аномалий. При этом идея packet classes из работы Bringhenti~\cite{bringhenti2020improving} перекликается с предложенным Diekmann разбиением пространства IP на классы эквивалентности и со структурой сервисных матриц, что подтверждает универсальность этой абстракции для отчётности.

\subsubsection{Анализ на основе журналов трафика}
\label{sec:Section2:logs}

Особняком стоит работа Golnabi, Min, Khan и Al-Shaer~\cite{golnabi2006analysis}, использующая методы data mining для анализа журналов межсетевого экрана. Авторы предлагают метод Mining firewall Log using Frequency (MLF) --- частотный анализ записей журнала, выводящий <<эффективные>> правила, отражающие реально наблюдаемый трафик, и метод Filtering-Rule Generalization (FRG) --- агрегацию примитивных правил в дерево обобщений по атрибутам пакета. Введены два дополнительных класса аномалий, не выявляемых статическим анализом: блокировка существующего сервиса и разрешение трафика к несуществующему сервису. Также вводится понятие \emph{decaying rules} (правил, чей трафик падает со временем) и \emph{dominant rules} (правил, обрабатывающих основную долю трафика).

Этот подход дополняет статический анализ, но не заменяет его. Для разрабатываемой системы данные журналов могут быть полезным дополнительным источником, однако опираться только на них опасно: отсутствие трафика, попадающего под правило, ещё не означает, что правило избыточно, --- сервис мог быть остановлен на время аудита. Поэтому log-анализ в настоящей работе рассматривается как направление дальнейшего развития инструментария, но не как ядро системы.

\subsection{Оптимизация политик: переупорядочивание и проектирование}
\label{sec:Section2:optimization}

Отдельный класс работ направлен не на обнаружение ошибок в существующей политике, а на её улучшение --- либо за счёт перестройки порядка правил, либо за счёт построения политики с нуля по высокоуровневой спецификации.

\subsubsection*{Переупорядочивание правил}

Значительный пласт работ посвящён задачам оптимизации производительности межсетевого экрана через переупорядочивание правил. Работа Fulp~\cite{fulp2005optimization} вводит DAG-модель precedence relations: для двух правил $R_i$ и $R_j$ с $i < j$ ребро $R_i \to R_j$ существует тогда и только тогда, когда их области пересекаются и действия различны. Любая топологическая сортировка такого DAG даёт семантически эквивалентный набор правил, что задаёт пространство допустимых перестановок и обосновывает корректность жадного алгоритма переупорядочивания по частоте срабатывания. Hamed и Al-Shaer~\cite{hamed2006dynamic} развивают эту линию и доказывают, что задача rule ordering optimization при наличии rule-dependency constraints \emph{NP-полна}; они предлагают двухслойную динамическую схему, где небольшое подмножество <<горячих>> правил отделяется от остальных и переупорядочивается по наблюдаемой частоте. Acharya, Wang, Ge, Znati и Greenberg~\cite{acharya2006traffic} дополняют эту линию явным разделением \emph{internal redundancy} (одна и та же область покрывается дублирующимися интервалами внутри одного правила) и \emph{external redundancy} (одно правило является супермножеством другого с тем же действием при меньшем ранге).

Эти работы являются классическими, однако их актуальность с точки зрения производительности снижается с развитием технологий аппаратного матчинга. На современных высоконагруженных сетевых устройствах (TCAM-чипы маршрутизаторов и коммутаторов, специализированные ASIC NGFW) поиск правила имеет асимптотику $O(1)$ либо логарифмическую, а не линейную, как было в эпоху чисто программных межсетевых экранов. Поэтому переупорядочивание правил с точки зрения производительности перестало быть критичным; гораздо большее значение имеет корректность набора. В разрабатываемой системе акцент делается именно на корректность; функция переупорядочивания может быть реализована как опциональная. Важнее то, что precedence DAG из работ Fulp и Hamed используется в разрабатываемом инструменте как формальный язык описания зависимостей между правилами для корректного предложения безопасных изменений администратору; формализация internal/external redundancy методически полезна и для аудита.

\subsubsection*{Проектирование политики с нуля и автоматизированный синтез}

Отдельный класс работ ориентирован на проектирование межсетевого экрана с нуля по высокоуровневой спецификации политики. Structured Firewall Design Gouda и Liu~\cite{gouda2007structured} (Computer Networks 2007) формализует три фундаментальные проблемы прямого проектирования рулсета как последовательности правил --- consistency, completeness и compactness --- и предлагает двухэтапный pipeline: сначала пользователь описывает политику в виде FDD, затем алгоритм автоматически конвертирует её в эквивалентный компактный список правил. Diverse Firewall Design Liu и Gouda~\cite{liu2008diverse} (IEEE TPDS 2008) переносит на firewall-проектирование принцип \emph{N-version programming}: одно требование выдаётся нескольким независимым командам, реализации сравниваются формальным алгоритмом (\emph{construction}, \emph{shaping}, \emph{comparison}) на базе ordered FDD, гарантированно находящим \emph{все} расхождения. Современная линия --- работы Bringhenti и др.~\cite{bringhenti2022automated} (IEEE TDSC 2023), формализующие задачу автоматической аллокации packet filter'ов и генерации их рулсетов в виртуализированном Service Graph как partial weighted MaxSMT-задачу. Hard-constraint'ами выступают сетевые требования безопасности (NSR), soft-constraint'ами --- минимизация числа firewall'ов и числа правил. Решение даёт \emph{correctness-by-construction}: итоговая политика гарантированно удовлетворяет всем NSR и оптимальна по выбранным метрикам.

В рамках того же направления стоит упомянуть Hybrid Tree-rule firewall Chomsiri, He, Nanda и Tan~\cite{chomsiri2016hybrid} (IEEE TCC 2016): правила задаются администратором через GUI в виде дерева (структурно близкого к FDD), что гарантирует \emph{конфликт-свободность по построению}; для runtime-обработки tree-rule конвертируется в listed rules с adaptive reordering по частоте срабатывания. При росте набора правил до 80\,000 потери производительности iptables составляют 47{,}66\%, тогда как hybrid tree-rule теряет лишь 7{,}43\%. Однако этот результат относится к программной обработке трафика; на современных hardware-firewall'ах с TCAM/ASIC преимущество tree-rule перед списком существенно сокращается.

В контексте настоящей работы следует отметить, что подобные подходы плохо учитывают реалии эксплуатации корпоративных сетей. Современные стандарты управления конфигурациями (Zero-Trust, ITSM-процессы, ISO/IEC 27001) предполагают привязку правил к заявкам пользователей, к управляемым объектам, к идентификаторам процессов и пр. Автоматический синтез политики <<с нуля>> разрушает эти связи и нарушает управляемость инфраструктуры. Кроме того, синтезированная политика может оказаться формально эквивалентной, но не отражающей принятые в организации стандарты. Diverse Firewall Design предполагает существование двух независимо разработанных версий, что в условиях реального заказчика крайне дорого; Structured Firewall Design требует, чтобы все правила выражались через интервалы целочисленных полей, что несовместимо с современными NGFW-расширениями; MaxSMT-подходы Bringhenti требуют формального описания NSR, отсутствующего в большинстве унаследованных конфигураций. Поэтому разрабатываемая в настоящей работе система не пытается порождать новые правила, а предоставляет администратору информацию о текущих ошибках для их направленного устранения. При этом отдельные технические идеи из синтезирующих работ --- канонический ordered FDD из Diverse Firewall Design, теоремы consistency/completeness FDD из Structured Firewall Design, представление политик через packet classes из работ Bringhenti --- используются как теоретические ориентиры для проектирования собственной модели правил.

Отметим также свежий обзор Bringhenti и др.~\cite{bringhenti2023automation} (ACM Computing Surveys 2023), систематизирующий современные подходы к автоматизации конфигурирования сетевой безопасности. Авторы предлагают декомпозицию задачи на две взаимосвязанные подзадачи: \emph{service composition} (выбор и аллокация network security functions) и \emph{function configuration} (генерация конфигурации каждой функции). Эта декомпозиция полезна для позиционирования настоящей работы: разрабатываемый инструмент относится исключительно к function configuration в режиме \emph{post-change аудита}, а не к synthesis или service composition. Из обзорной литературы также заслуживает упоминания работа Ramesh и др.~\cite{ramesh2024exploring} (Indiana Journal of Multidisciplinary Research, 2024), систематизирующая направления rule prioritization, complexity management, anomaly detection, optimization и роль AI/ML в управлении политиками; используемая в ней таксономия отношений правил (IM/EM/PM/CD/C) согласуется с предложенной Al-Shaer и Hamed и применяется в настоящей работе.

\subsection{Сводка существующих подходов и обоснование собственной модели}
\label{sec:Section2:summary}

Краткая сводная таблица рассмотренных подходов приведена в таблице~\ref{tab:methods}. В первой колонке указано основное направление работы, во второй --- ключевая структура данных или формализм, в третьей --- область применимости, в четвёртой --- основные ограничения с точки зрения настоящей работы.

\begin{table}[h!]
\centering
\caption{Сводка основных подходов к анализу правил фильтрации}
\label{tab:methods}
\begin{tabular}{|p{3.5cm}|p{3.5cm}|p{3.5cm}|p{4.0cm}|}
\hline
\textbf{Подход} & \textbf{Структура данных} & \textbf{Область} & \textbf{Ограничения} \\
\hline
Эмпирический\,~\cite{wool2004quantitative,wool2010trends} & Списки ошибок, метрики сложности & Аудит реальных МЭ & Нет автоматического алгоритма \\
\hline
Pairwise (Al-Shaer)\,~\cite{al2003firewall} & Policy tree, попарные предикаты & 4 класса аномалий + editor & Не находит групповые \\
\hline
Геометрия конфликтов\,~\cite{hari2000detecting,eppstein2000internet} & Прямоугольники в $d$-простр., resolve-фильтры & Конфликты пар правил, $O(n^{3/2})$ & Только пары \\
\hline
Query-based (Fang)\,~\cite{mayer2000fang} & Topology + zones + simulation & Multi-vendor query engine & Lightweight testing; не exhaustive \\
\hline
FIREMAN\,~\cite{yuan2006fireman} & BDD, состояния $\langle A,D,F\rangle$ & Полный анализ конфигурации & Только Cisco PIX/PF; нет NAT \\
\hline
FDD\,~\cite{gouda2004firewall,liu2008formal} & Acyclic graph по полям & Проектирование, верификация & Сложен импорт реальных правил \\
\hline
Acharya--Gouda\,~\cite{acharya2011firewall} & Projection firewall & Verif. $\Leftrightarrow$ Redundancy & Stateless; одиночное правило \\
\hline
Diekmann\,~\cite{diekmann2018verified} & Isabelle/HOL, ternary logic & Полный анализ iptables & Только iptables-filter \\
\hline
SAT\,~\cite{zhang2012verification} & CNF, Tseitin & Eq./Inclusion checking & Synthesis непрактичен \\
\hline
Margrave\,~\cite{nelson2010margrave} & FOL + Kodkod (SAT) & IOS, scenario-finding & Только Cisco IOS \\
\hline
Change-impact\,~\cite{liu2008firewall} & FDD + diff & Pre-change анализ & Не для post-change аудита \\
\hline
TLA$^+$\,~\cite{kapus2023improved} & TLC model checker & SDN/OpenFlow сети & Только SDN \\
\hline
HSA\,~\cite{kazemian2012header,bringhenti2020improving} & Header space, packet classes & Reachability/loops/isolation & Не находит intra-FW аномалий \\
\hline
Log-mining\,~\cite{golnabi2006analysis} & ARM, частотный анализ & Decaying rules & Дополнение, не основа \\
\hline
Reordering\,~\cite{fulp2005optimization,hamed2006dynamic,acharya2006traffic} & DAG precedence, NP-полнота & Производительность & Неактуально для O(1) hardware \\
\hline
Tree-rule\,~\cite{chomsiri2016hybrid} & Дерево правил, GUI & Конфликт-свобода by design & Нужна перестройка с нуля \\
\hline
Synthesis\,~\cite{liu2008diverse,gouda2007structured,bringhenti2022automated} & FDD, MaxSMT & Проектирование с нуля & Разрушает учётность правил \\
\hline
\end{tabular}
\end{table}

Из проведённого анализа можно сформулировать следующие требования к разрабатываемой системе.

\textbf{Унифицированная вендоро-независимая модель}. Модель должна позволять представлять правила различных платформ (классические ACL Cisco/Juniper/Huawei, отечественные NGFW, iptables) единообразно. На уровне атрибутов это означает поддержку как примитивных квалификаторов (адреса, порты, протоколы), так и именованных объектов и групп. На уровне действий --- базовое разделение на \emph{allow}/\emph{deny} с возможностью различения дополнительных классов действий и состояния \emph{Undecided} в смысле Diekmann для chain-семантик.

\textbf{Обработка неизвестных match-условий}. Опираясь на работу Diekmann, в систему должна быть встроена техника upper/lower closure для match-условий, которые невозможно интерпретировать на уровне модели. При этом аппроксимация должна быть информативной: вместо тотальной потери информации, для неизвестных match-условий следует использовать пятёрку из правила как контекст, остальные поля считая допускающими произвольное значение.

\textbf{Поиск интрафайервольных аномалий полного покрытия}. По образцу FIREMAN система должна оперировать состоянием $\langle A,D,F\rangle$ или его эквивалентом и обнаруживать классы shadowing, redundancy, generalization, correlation, а также irrelevance. Union-аномалии могут быть обнаружены, но представляются предупреждением, а не ошибкой, ввиду указанной выше специфики управления правилами в Zero-Trust-инфраструктурах.

\textbf{Опора на методологию статического post-change аудита}. Разрабатываемая система предназначена для регулярной проверки уже эксплуатируемых конфигураций, а не для анализа предполагаемых изменений. Это определяет приоритеты: точность и полнота обнаружения аномалий важнее, чем интерактивная производительность; объём отчёта должен быть достаточным для администратора, но не разрозненным; результат должен быть привязан к конкретным правилам исходной конфигурации.

\textbf{Двухфазная архитектура нормализатор + анализатор}. Парсинг и нормализация исходных конфигураций отделяются от ядра анализа в самостоятельный модуль. Это обеспечивает расширяемость на новые форматы и позволяет вводить поддержку отечественных платформ инкрементно, без переработки ядра.

\textbf{Возможность расширения}. Архитектура системы должна оставлять место для функций, не входящих в основной набор: change-impact-анализ, log-mining, reachability-анализ, интеграция с системами учёта правил.

С учётом перечисленных требований основными ориентирами для архитектуры разрабатываемого инструмента являются работы Yuan и др.~\cite{yuan2006fireman} и Diekmann и др.~\cite{diekmann2018verified}. От FIREMAN заимствуются формальная модель состояний $\langle A,D,F\rangle$ при анализе ACL, классификация аномалий и структура отчёта. От Diekmann --- двухфазная архитектура нормализатор + анализатор, тернарная логика для неизвестных match-условий, идея сервисных матриц как формы сводного отчёта. Дополнительно используется унификация представления правил, описанная в обзоре Ramesh и др.~\cite{ramesh2024exploring} (отношения IM/EM/PM/CD/C). Совокупность этих идей служит основанием для разрабатываемой в последующих разделах унифицированной вендоро-независимой модели правил и алгоритмов анализа на её основе.

\newpage
